\gddchapter{Interview sequence}


Below is the guideline for the semi-open guided interview. 
The interview starts with easy warm-up questions and progresses towards harder questions as it goes along. 

The answers were translated from polish back to english.

\section{Introduction}
Remind the participant that this is not a test, rather a joint exploration.

\section{Icebreakers}

\begin{QandA}
   \item Now that you’ve tried both versions, how are you feeling? Was there anything that immediately surprised you?
         \begin{answered}
               well using the second version when I could just pick text seems much easier but gives us less opportunities to try something else so I would rather if I was to play this game for longer I would rather go for the first version so I would rather speak to the game talk to the game and just wonder if it's
               understands me or not, or just even try several times, but I think it's more interesting, more challenging.
               So yeah, this is my opinion, if this was the question.
               Although, the second version makes it easier to navigate, because you can already see the options, the opportunities.
               and in the first option I have to try several times before I pick the right one, possible one.
            \end{answered}

   \item If you had to describe the difference between these two games to a person who hasn't played either, what would you say?
         \begin{answered}
               well using the second version when I could just pick text seems much easier but gives us less opportunities to try something else so I would rather if I was to play this game for longer I would rather go for the first version so I would rather speak to the game talk to the game and just wonder if it's
               understands me or not, or just even try several times, but I think it's more interesting, more challenging.
               So yeah, this is my opinion, if this was the question.
               Although, the second version makes it easier to navigate, because you can already see the options, the opportunities.
                nd in the first option I have to try several times before I pick the right one, possible one.


               Gabriel: And if you had to describe the difference between those two games to a person who haven't played at all, how would you describe it?
               I would say that, okay, that the first version reminds me a bit of those games from the 80s or 90s.
                kind of games from my childhood that seems like everyone knows what I'm talking about and everyone knows that this is kind of a nice thing to do and just picking the, not writing, but picking the right, the options you want to choose.
               I just, once again, what was the question?
               Because I forgot.


               Gabriel: If you had to describe the difference between two versions, to a person who hasn't played either, what would you say to them?
           
               Klara: So the main difference is that the first version makes me think more of a chat GPT stuff and AI conversation and the other version makes me think of those old 80s games from Sierra or something.
               This kind of, you know, first computer games that appeared in 1990s, which actually means a lot of me because I really like playing them.
               So this will be my main, like this first version is more
               up-to-date because it makes me think it's more connected to the AI and the fact that you can talk to the computer and it understands you and the other version is much, let's say, simpler and you just pick one of four options and the game perfectly knows what to do.
               So the first version gives you more of an unknown result, while the second version is more predictable, I would say.
               Yeah, that would be, I think, it.
         \end{answered}
\end{QandA}




\section{Grand tour questions}
(remember to pause here 3-5 seconds)

\begin{QandA}


   \item From the moment you first used your voice to navigate, what were you thinking as you decided what to say to the system?

         \begin{answered}
            Klara: How to take the first step?
             How to start?
            That was my main concern.
            If I understand properly the situation I'm in and where should I start from?
            To be... Maybe not to make a mistake, because I'm sure you can't make a mistake while playing, but just to be more effective and just to head towards the goal, let's say.
            So what should I do first?
            Let's move.
            My idea for the game was actually to move a lot, so I was thinking how to take the first step.


            Gabriel: I'm trying to also understand what made you choose the words that you chose.
            What I mean by that is that, for instance, some people would only speak in first person because they have their reasons, while other people, for instance, decide to only speak to the computer as if it's an object. or a person in the room.
            So what made you make your choice and how did you use that?


            Klara: Yes, I never even thought about it.
            It was just natural.
            let's go here, let's go there, I want to do this, I want to do that, that was obvious to me, I mean, I feel like we are playing together, me and the computer, and not like I'm asking him to take me somewhere, it's like we are together in this, and yeah, that's why I decided, I even, I think I even used this form, let's go here, so I mean like we're playing together, like me and him, like both of us, so it came naturally
         \end{answered}

    \item How was it like to use the voice input in the game?
         \begin{answered}
            Klara: Quick and easy to use.
            You mean the voice input, so I don't have to...
            I can say more than just picking the lines that are ready and they're waiting for me.
            So actually I can be more specific and be more challenging to the game, I guess.
            I guess I can ask the game for more than it maybe understands.


            Gabriel: How did it feel like?


            Klara: Good.
            I'm quite used to talk to the computer since the AI spreads around.
            getting more and more normal to talk to machines, so that's fine.
            Didn't have any problem with that.
            The reason why I was just wondering if it's gonna... if it understands me while I'm wearing a mask, but I think it went well.
         \end{answered}


\end{QandA}




\section{Core comparative questions}

\begin{QandA}
   \item In which version did you feel more 'inside' the story world?
         \begin{answered}
            Klara: in the second version which can be it's tricky but yeah in the second version because maybe it's also maybe because I already got to know the game a little bit because the first version was like I didn't know anything about the game and I had to get used to this and to reading a lot because there's a lot to be read first a very nice description but still it builds this atmosphere and you have to read it carefully and I didn't read it till the end and that was the mistake
            yeah but this this means that the first the beginning was a little bit I had to get too used to that but so the second part was easier in fact for me because not only because there were already options waiting for me but also because I already knew the game a little bit I knew what to expect and I knew this mechanic of the game I know I had to write the description till the very end because at the very end there are
            tips, what I can do, hints, let's say.


            Gabriel: Just to spend a little bit more time here.
            I understand the difficulty, right?
            But at the same time, I'm trying to probe here and ask, which of those versions felt more like, how would I explain this?
            In which of those did you feel more convinced that you are actually
            in the game, let's say. Kind of like when you watch a movie sometimes, and the movie makes you forget that you... makes you feel like, you know, you are in the movie, makes you forget about the fact that you're actively watching it.


            Klara: The first version, I think.
            The first version, because... because... because...
            I actually don't know how to justify that, but this is just a gut feeling that the first version
            

           
         \end{answered}

   \item How did the 'effort' of thinking of what to say naturally compare to the 'effort' of knowing which keys to press?

         \begin{answered}
            oh it's [the voice version] much bigger effort because i have to i have to first of all you need to read carefully you need to be patient because in the text there are hints and places described and places to go and the directions you can take
            While the second version is actually can mostly skip the text if you don't like reading I mean I like reading because as I mentioned before the text was really Interesting and very well Written but if someone does not like reading too much or be just in a hurry and wants to play fast I would say the second version is easier because you know what to
            you can do.
            You know what the game's limits are.
            And in the first version, I have no idea about the limits.
            I don't know if it's going to put me on the moon if I want to or not.
            In the second version, I know that I only have four options to choose from, and that's mostly that.
         \end{answered}

    \item Did the freedom to say 'anything' make you feel more in control, or was it more overwhelming than having a fixed list of keys?

         \begin{answered}
            more overwhelming, but also because this was my first experience with the game.
            Now if I was to play another time, I would probably be more focused on really choosing the plan, making a plan for the game and not really trying to like, you know, coordinate somehow and be understood by the game.
            Now I would be, a second take would be, for me, I think much easier.
         \end{answered}
    \item How would you compare the overall experience of using the two game versions? Which one would you choose and why?

         \begin{answered}
            Klara: Yeah, the first version, I guess, the first version.
            But then I would concentrate more on reading texts and see where the game...
            I mean, I would like to check the game's limits.
            That would be my... the biggest fun for me.
            And I suppose that the game creator could be also very inventive and hide some crazy solutions for those people who want to play this way and say, okay, let's just, I don't know, go to the roof and then you go to the roof of the building and like to be more...
            give more place and space to explore for those people who want to talk to the game and not just pick the already answered out there already.


            Gabriel: Actually, it can be done in this game.
            You can go to the room.
            Yeah, good.
            But of course, of course.


            Klara: It takes time probably, you have to play some time.


            Gabriel: Yeah, that's absolutely right.


            Klara: You'll come to that, yeah.


            Klara: But I fully feel you're right, having the ability to test, stress test, kind of, right, see what can I do.
            Can I go to the moon or not?
            And where are the surprises and presents for me?
            Because I expect some surprises and presents in this version for heart-seeking players, players who are patient and want to explore more than just go through the game quickly and


            Gabriel: Yeah, you like, in some video games there are these things like hidden secrets, right?


            Klara: Exactly.
            This is it.
            This is what I would expect from this first version.
            It gives me more opportunities and then I can really, okay, I can even ask him to go and pick this, this, this and check a few places around the room even.
            Not just to go up and down but let's say let's check this box, let's check this box, let's, I don't know,
            see if there's anything on the floor, this kind of things.
            

            
         \end{answered}
\end{QandA}






\section{Probing}
In this part of the interview look into the sticky moments noted above and ask follow-up questions probing for more detail.
Include lexical reflection (eg. you used the word "clunky" when explaining navigation. Could you tell me more about what was happening
in that moment?)

The probing was done as follow up questions to the pre-existing question list rather than separate questions in this test.


\section{Final reflection}

\begin{QandA}
   \item After reflecting on everything today, is there anything else you’d like to add about using your voice to play this game that we haven't talked about?

         \begin{answered}
            Oh, yeah, I think I very much like the descriptions.
            They are absolutely amazing and very like...
            very romantic I'd say even, perfect description, so I'd like to emphasize that.
            I would like to see, if I was to play that longer, I would like to see a bit more of the picture because it made me somehow suffocate that I could not see the whole picture and I could only see the picture, half of the picture in fact.
            Well, this was my impression.
            So if I was to change something to the game, add something from myself, let's say, a general comment, I would like to see, because I really like the pictures as well, the images, I would like to see more of the image and be able to pick more details on it.
            Because each of the images were really creative and we'd like to see more, more of this.
            And I think the font should be a bit bigger because it's very, very small and hard to read on this kind of
            nice but small screen anyway.
         \end{answered}
\end{QandA}

\section{Closing}
Thank the participant for his time and tell them again how much value his input and time gives me.






